\hmsubsection{Data Flow}{FS}{}\label{sec:data-flow}
This section describes what happens during one pick-and-place cycle,
step by step. The cycle starts when the supervisor enters the
\texttt{SCAN} state with the arm at \texttt{SCAN\_POSE}. It ends when
the arm returns to \texttt{SCAN\_POSE} after dropping a ball in the
right bin.
\subsubsection*{Step 1 -- Frame acquisition}
The OAK-D S2 takes a colour frame at 640 by 400 pixels and sends it
to the Pi over USB 3.0. The DepthAI driver delivers the frame as a
NumPy array in BGR format. One frame per loop iteration, at roughly
20--25 Hz.
\subsubsection*{Step 2 -- Perception}
The frame goes through the perception layer. The output is a list of
tracked detections. Each detection has a colour, a pixel centre, a
radius, a confidence, and a stable ID.
\subsubsection*{Step 3 -- Target selection}
If both red and blue balls are visible, the supervisor picks one based
on a priority rule. The default rule is red first, blue second. If no
detection is present for two frames in a row, the supervisor returns
to \texttt{IDLE} without moving the arm.
\subsubsection*{Step 4 -- Pixel to workspace}
The pixel centre of the chosen detection is converted to workspace
centimetres through the perspective transformation from the touch
calibration (Section \ref{sec:touch-cal}). The result is a target
$(x, y, z)$ in the workspace frame. The $z$ value comes from the
per-point surface height recorded during touch calibration. For
positions between calibration points, the height is interpolated.
\subsubsection*{Step 5 -- Inverse kinematics}
The target goes into \texttt{ArmIK} (Section \ref{sec:ik-solver}). The
solver returns five Dynamixel goal positions in encoder steps. If the
target is out of reach, the solver either clamps it or raises an
exception, depending on the operating mode
(Section \ref{sec:joint-limits}).
\subsubsection*{Step 6 -- Trajectory generation}
The supervisor turns the single goal pose into three waypoints: an
approach pose 10 cm above the target, the pick pose at the target,
and a lift pose 10 cm above the target. Each waypoint is solved
independently and queued for transmission.
\subsubsection*{Step 7 -- Command transmission}
The communication layer sends each motor command over the USB serial
link to the OpenRB-150 as an ASCII line. The firmware replies with
\texttt{ACK}, then \texttt{BUSY} during motion, then \texttt{READY}
when the encoders settle within the tolerance band.
\subsubsection*{Step 8 -- Gripper actuation}
At the pick pose the supervisor closes the gripper. At the release
pose over the right bin, the supervisor opens it. The gripper uses
the same protocol; the firmware translates \texttt{GRIP close} and
\texttt{GRIP open} into the pre-calibrated goal positions.
\subsubsection*{Step 9 -- Return to scan}
After the ball is released, the supervisor sends the arm back to
\texttt{SCAN\_POSE}. Once the arm settles, the supervisor returns to
the \texttt{SCAN} state. No scan-pose current limit is applied in the
deployed system, since the tested limit did not sufficiently reduce
the elbow motor's thermal load (Section \ref{sec:thermal}). Instead,
the deployed scan pose relies on the mechanically supported posture in
which the L\textsubscript{2} structure rests on the L\textsubscript{1}
piece while preserving the camera's view of the workspace. This reduces
idle holding torque at \texttt{SCAN\_POSE}, but movement-related motor
heating remains outside this scan-state mitigation.
\subsubsection*{Step 10 -- Logging}
One row is appended to the cycle log. The row has a timestamp, the
colour of the ball, the pixel and workspace coordinates, the per-stage
timing (acquisition, perception, planning, motion), and an outcome
flag. The log is the main source of evidence for the tests described
in Chapter \ref{cha:testing}.
